%% kapa: exam=1 round=32 year=2021 period=1 form=A subject=civil_law subject_ko=민법 seq=34 number=34
%% source: Q-Net 감정평가사 기출문제 (cst00309) artlSeq=5208397 file=extracted/1008/5208397/A형/2021년도 제32회 감정평가사 1차 1교시 A형.pdf
%% conversion: kordoc 4.24.51 -> tools/gen1.py (auto); review=proofread-sonnet; flags=none
\begin{question}{34}
전세권에 관한 설명으로 옳은 것은? (다툼이 있으면 판례에 따름)
\begin{choices}
  \item 전세권이 성립한 후 목적물의 소유권이 이전되더라도 전세금반환채무가 당연히 신소유자에게 이전되는 것은 아니다.
  \item 전세권의 존속기간이 시작되기 전에 마친 전세권설정등기는 특별한 사정이 없는 한 그 기간이 시작되기 전에는 무효이다.
  \item 전세권을 설정하는 때에는 전세금이 반드시 현실적으로 수수되어야 한다.
  \item 건물의 일부에 전세권이 설정된 경우 전세권의 목적물이 아닌 나머지 부분에 대해서도 경매를 신청할 수 있다.
  \item 전세권자가 통상의 필요비를 지출한 경우 그 비용의 상환을 청구하지 못한다.
\end{choices}
\end{question}
\answer{5}
